\chapter{Aritmética}\label{ch:g12:arith}

A aritmética estuda os inteiros: \omterm{def:g12:arith:divides}{divisibilidade}, \omterm{def:g12:arith:prime}{números primos}, restos.
Durante muito tempo considerada a mais pura da matemática pura, ela hoje
protege cada pagamento on-line: o criptossistema RSA repousa sobre os
teoremas de Bézout, Gauss e Fermat demonstrados neste capítulo.

\section{Divisibilidade e divisão euclidiana}

\begin{definition}[Divisibilidade]\label{def:g12:arith:divides}
Sejam $a, b \in \Z$. Dizemos que $b$ \emph{divide}
$a$\index{divisibilidade}, o que se escreve $b \mid a$, se existe
$k \in \Z$ com $a = kb$. Dizemos também que $a$ é \emph{múltiplo} de $b$.
\end{definition}

\begin{proposition}\label{prop:g12:arith:divprops}
Se $c \mid a$ e $c \mid b$, então $c$ \omterm{def:g12:arith:divides}{divide} toda combinação inteira
$au + bv$ ($u, v \in \Z$). Se $a \mid b$ e $b \mid a$ com $a,b \in \N$,
então $a = b$. Se $a \mid b$ e $b \neq 0$, então $\abs a \leq \abs b$.
\end{proposition}

\begin{proof}
Escreva $a = kc$, $b = lc$: então $au + bv = (ku + lv)c$. Os demais itens
seguem de $\abs{a} = \abs{k}\,\abs{b}$ com $\abs k \geq 1$ quando
$b = ka \neq 0$.
\end{proof}

\begin{theorem}[Divisão euclidiana]\label{thm:g12:arith:euclid}
\index{divisão euclidiana}
Sejam $a \in \Z$ e $b \in \N^*$. Existe um único par
$(q, r) \in \Z \times \N$ tal que
\[
a = bq + r \qquad\text{e}\qquad 0 \leq r < b .
\]
$q$ é o \emph{quociente} e $r$ é o \emph{resto}.
\end{theorem}

\begin{proof}
\emph{Existência.} O conjunto dos múltiplos de $b$ que não excedem $a$ tem
um maior elemento $bq$ (ele é não vazio e \omterm{def:g12:seq:bounded}{limitado} superiormente); ponha
$r = a - bq$. Pela maximalidade, $b(q+1) > a$, logo $0 \leq r < b$.
\emph{Unicidade.} Se $bq + r = bq' + r'$ com $0 \leq r, r' < b$, então
$b(q - q') = r' - r$ e $\abs{r' - r} < b$: um múltiplo de $b$ de valor
absoluto menor que $b$ só pode ser $0$, logo $r = r'$ e $q = q'$.
\end{proof}

\section{Congruências}

\begin{definition}[Congruência]\label{def:g12:arith:congruence}
Seja $n \in \N^*$. Dois inteiros $a, b$ são \emph{congruentes módulo
$n$}\index{congruência}, o que se escreve $a \equiv b \pmod n$, se
$n \mid (a - b)$ --- equivalentemente, se $a$ e $b$ têm o mesmo resto na
\omterm{thm:g12:arith:euclid}{divisão euclidiana} por $n$.
\end{definition}

\begin{proposition}[Compatibilidade com as operações]\label{prop:g12:arith:congops}
Se $a \equiv b \pmod n$ e $c \equiv d \pmod n$, então
\[
a + c \equiv b + d, \qquad
ac \equiv bd, \qquad
a^k \equiv b^k \ (k \in \N) \pmod n .
\]
\end{proposition}

\begin{proof}
$n$ \omterm{def:g12:arith:divides}{divide} $(a-b) + (c-d) = (a+c) - (b+d)$, e
$ac - bd = a(c - d) + d(a - b)$ também é múltiplo de $n$. A regra das
potências segue por indução a partir da regra do produto.
\end{proof}

\begin{method}[Calcular potências módulo $n$]\label{met:g12:arith:powers}
Para calcular $a^k \bmod n$, reduza a base módulo $n$, depois procure uma
potência pequena de $a$ congruente a $\pm1$ e use-a para encolher o
expoente. Por exemplo, $2^{100} \bmod 7$: como
$2^3 = 8 \equiv 1 \pmod 7$ e $100 = 3\times33 + 1$,
\[
2^{100} = \left(2^{3}\right)^{33} \times 2 \equiv 1^{33}\times 2 = 2 \pmod 7 .
\]
\end{method}

\section{MDC, Bézout e Gauss}

\begin{definition}[MDC]\label{def:g12:arith:gcd}
Sejam $a, b$ inteiros não ambos nulos. O \emph{máximo divisor
comum}\index{mdc} $\gcd(a, b)$ é o maior inteiro que \omterm{def:g12:arith:divides}{divide} $a$ e $b$.
Quando $\gcd(a,b) = 1$, diz-se que $a$ e $b$ são \emph{primos entre
si}\index{primos entre si}.
\end{definition}

\begin{proposition}[Algoritmo de Euclides]\label{prop:g12:arith:euclidalgo}
\index{algoritmo de Euclides}
Se $a = bq + r$ ($b \neq 0$), então $\gcd(a, b) = \gcd(b, r)$. Iterar a
\omterm{thm:g12:arith:euclid}{divisão euclidiana}, portanto, calcula $\gcd(a,b)$: o \omterm{def:g12:arith:gcd}{mdc} é o último resto
não nulo.
\end{proposition}

\begin{proof}
Todo divisor comum de $a$ e $b$ \omterm{def:g12:arith:divides}{divide} $r = a - bq$
(\cref{prop:g12:arith:divprops}) e é, portanto, divisor comum de $b$ e
$r$; e reciprocamente, já que $a = bq + r$. Os dois pares têm os mesmos
divisores comuns e, portanto, o mesmo \omterm{def:g12:arith:gcd}{mdc}. O algoritmo termina porque os
restos formam uma \omterm{def:g12:seq:sequence}{sequência} estritamente \omterm{def:g10:functions:variations}{decrescente} de inteiros não
negativos.
\end{proof}

\begin{example}\label{ex:g12:arith:euclidalgo}
$\gcd(252, 198)$: $252 = 198 + 54$; $198 = 3\times54 + 36$;
$54 = 36 + 18$; $36 = 2 \times 18 + 0$. Logo
$\gcd(252,198) = 18$.
\end{example}

\begin{theorem}[Identidade de Bézout]\label{thm:g12:arith:bezout}
\index{identidade de Bézout}
Sejam $a, b$ inteiros não ambos nulos e $d = \gcd(a,b)$. Existem
$u, v \in \Z$ tais que
\[
au + bv = d .
\]
Em particular, $a$ e $b$ são \omterm{def:g12:arith:gcd}{primos entre si} se, e somente se,
$au + bv = 1$ para alguns inteiros $u, v$.
\end{theorem}

\begin{proof}
Percorra o \omterm{prop:g12:arith:euclidalgo}{algoritmo de Euclides} de trás para a frente: cada resto é
combinação inteira dos dois anteriores, e os dados iniciais $a, b$ são
combinações de si mesmos; por substituição descendente, o último resto não
nulo $d$ é combinação inteira de $a$ e $b$. (No
\cref{ex:g12:arith:euclidalgo}:
$18 = 54 - 36 = 54 - (198 - 3\times54) = 4\times54 - 198 =
4(252 - 198) - 198 = 4\times252 - 5\times198$.)

Para a equivalência: se $\gcd(a,b) = 1$, Bézout fornece $u, v$;
reciprocamente, todo divisor comum de $a$ e $b$ \omterm{def:g12:arith:divides}{divide} $au + bv = 1$, o
que obriga $\gcd(a,b) = 1$.
\end{proof}

\begin{theorem}[Lema de Gauss]\label{thm:g12:arith:gauss}
\index{lema de Gauss}
Sejam $a, b, c \in \Z$. Se $a \mid bc$ e $\gcd(a, b) = 1$, então
$a \mid c$.
\end{theorem}

\begin{proof}
Bézout dá $au + bv = 1$; multiplique por $c$: $acu + bcv = c$. Os dois
termos do lado esquerdo são múltiplos de $a$ (o segundo porque
$a \mid bc$), logo $c$ também é.
\end{proof}

\begin{corollary}\label{cor:g12:arith:coprimeprod}
Se $a \mid c$, $b \mid c$ e $\gcd(a,b) = 1$, então $ab \mid c$.
\end{corollary}

\begin{proof}
Escreva $c = ak$. De $b \mid ak$ e $\gcd(a,b)=1$, Gauss dá $b \mid k$,
digamos $k = bl$; então $c = abl$.
\end{proof}

\section{Números primos}

\begin{definition}[Primo]\label{def:g12:arith:prime}
Um inteiro $p \geq 2$ é \emph{primo}\index{número primo} se seus únicos
divisores positivos são $1$ e $p$.
\end{definition}

\begin{proposition}\label{prop:g12:arith:primedivides}
Todo inteiro $n \geq 2$ tem um divisor \omterm{def:g12:arith:prime}{primo}; se $n$ não é \omterm{def:g12:arith:prime}{primo}, ele tem
um divisor \omterm{def:g12:arith:prime}{primo} $\leq \sqrt n$. Se um \omterm{def:g12:arith:prime}{primo} $p$ \omterm{def:g12:arith:divides}{divide} um produto $ab$,
então $p \mid a$ ou $p \mid b$ (\emph{lema de Euclides}).
\end{proposition}

\begin{proof}
O menor divisor $d \geq 2$ de $n$ é \omterm{def:g12:arith:prime}{primo} (qualquer divisor próprio de $d$
seria um divisor menor de $n$). Se $n = de$ é composto com
$2 \leq d \leq e$, então $d^2 \leq de = n$, logo $d \leq \sqrt n$. Para o
lema de Euclides: se $p \nmid a$, então $\gcd(p, a) = 1$ (os únicos
divisores de $p$ são $1$ e $p$), e o lema de Gauss dá $p \mid b$.
\end{proof}

\begin{theorem}[Euclides]\label{thm:g12:arith:infiniteprimes}
Há infinitos \omterm{def:g12:arith:prime}{números primos}.
\end{theorem}

\begin{proof}
Dada uma lista finita qualquer $p_1, \dots, p_k$ de \omterm{def:g12:arith:prime}{primos}, considere
$N = p_1 p_2 \cdots p_k + 1$. Algum \omterm{def:g12:arith:prime}{primo} $p$ \omterm{def:g12:arith:divides}{divide} $N$; mas nenhum $p_i$
\omterm{def:g12:arith:divides}{divide} $N$ (o resto é $1$), de modo que $p$ é um \omterm{def:g12:arith:prime}{primo} fora da lista.
Nenhuma lista finita esgota os \omterm{def:g12:arith:prime}{primos}.
\end{proof}

\begin{theorem}[Teorema fundamental da aritmética]\label{thm:g12:arith:fta}
\index{teorema fundamental da aritmética}
Todo inteiro $n \geq 2$ é produto de \omterm{def:g12:arith:prime}{primos}, e essa fatoração é única a
menos da ordem dos fatores:
\[
n = p_1^{\alpha_1} p_2^{\alpha_2} \cdots p_r^{\alpha_r},
\qquad p_1 < p_2 < \dots < p_r \text{ primos},\ \alpha_i \geq 1 .
\]
\end{theorem}

\begin{proof}
\emph{Existência}, por indução forte: se $n$ é \omterm{def:g12:arith:prime}{primo}, ele é sua própria
fatoração; caso contrário, $n = de$ com $2 \leq d, e < n$, e ambos se
fatoram pela hipótese de indução. \emph{Unicidade}: suponha
$p_1\cdots p_s = q_1 \cdots q_t$ (\omterm{def:g12:arith:prime}{primos}, com repetições permitidas). Pelo
lema de Euclides, $p_1$ \omterm{def:g12:arith:divides}{divide} algum $q_j$ e, sendo \omterm{def:g12:arith:prime}{primo}, $p_1 = q_j$;
cancele e repita. As duas fatorações coincidem termo a termo.
\end{proof}

\begin{theorem}[Pequeno teorema de Fermat]\label{thm:g12:arith:fermat}
\index{pequeno teorema de Fermat}
Seja $p$ \omterm{def:g12:arith:prime}{primo} e $a \in \Z$ com $p \nmid a$. Então
\[
a^{p-1} \equiv 1 \pmod p .
\]
Para todo $a \in \Z$ (sem hipótese de coprimalidade),
$a^p \equiv a \pmod p$.
\end{theorem}

\begin{proof}
Considere os $p - 1$ inteiros $a, 2a, 3a, \dots, (p-1)a$ módulo $p$.
Nenhum é $\equiv 0$ (se $p \mid ka$ com $1 \leq k \leq p-1$, o lema de
Euclides obriga $p \mid k$, o que é impossível), e eles são dois a dois
distintos módulo $p$ (se $ka \equiv la$, então $p \mid (k - l)a$, logo
$p \mid k - l$, logo $k = l$). Assim, módulo $p$, eles são os números
$1, 2, \dots, p-1$ em alguma ordem. Multiplicando todas as \omterm{def:g12:arith:congruence}{congruências}:
\[
a^{p-1}\,(p-1)! \equiv (p-1)! \pmod p .
\]
Como $p$ não \omterm{def:g12:arith:divides}{divide} nenhum entre $1, \dots, p-1$, usar repetidamente o
lema de Euclides permite cancelar $(p-1)!$, restando
$a^{p-1} \equiv 1$. A segunda forma segue multiplicando por $a$ (e é
trivial quando $p \mid a$).
\end{proof}

\begin{example}[Aplicação à criptografia]\label{ex:g12:arith:rsa}
O teorema de Fermat torna reversível a exponenciação módulo $n$ quando os
expoentes são bem escolhidos --- o coração do criptossistema \emph{RSA}.
Com $p, q$ \omterm{def:g12:arith:prime}{primos} grandes e $n = pq$, publicam-se $n$ e um expoente $e$; a
cifragem é $x \mapsto x^e \bmod n$. Decifrar exige um expoente $d$ com
$ed \equiv 1 \pmod{(p-1)(q-1)}$, que só quem conhece $p$ e $q$ consegue
calcular --- e recuperar $p, q$ a partir de $n$ significa \omterm{def:g10:algebra:expand}{fatorar} um número
de centenas de algarismos, o que nenhum algoritmo conhecido faz em tempo
razoável.
\end{example}

\section{Exercícios}

\begin{exercise}[$\star$]\label{exo:g12:arith:1}
Calcule o quociente e o resto da \omterm{thm:g12:arith:euclid}{divisão euclidiana} de $2026$ por $17$ e
de $-2026$ por $17$.
\end{exercise}

\begin{exercise}[$\star$]\label{exo:g12:arith:2}
Qual é o resto de $7^{100}$ módulo $10$? (Qual é o último algarismo de
$7^{100}$?)
\end{exercise}

\begin{exercise}[$\star$]\label{exo:g12:arith:3}
Usando o \omterm{prop:g12:arith:euclidalgo}{algoritmo de Euclides}, calcule $\gcd(1071, 462)$ e encontre
inteiros $u, v$ com $1071u + 462v = \gcd(1071, 462)$.
\end{exercise}

\begin{exercise}[$\star$]\label{exo:g12:arith:4}
Mostre que, para todo $n \in \Z$, $n^2$ é congruente a $0$ ou $1$ módulo
$4$. Deduza que um inteiro $\equiv 3 \pmod 4$ nunca é soma de dois
quadrados.
\end{exercise}

\begin{exercise}[$\star\star$]\label{exo:g12:arith:5}
Mostre que, para todo $n \in \N$, $n(n+1)(2n+1)$ é divisível por $6$.
\end{exercise}

\begin{exercise}[$\star\star$]\label{exo:g12:arith:6}
Resolva em $\Z$ a \omterm{def:g12:arith:congruence}{congruência} $5x \equiv 3 \pmod{11}$. (Sugestão:
encontre o inverso de $5$ módulo $11$.)
\end{exercise}

\begin{exercise}[$\star\star$]\label{exo:g12:arith:7}
Resolva em $\Z \times \Z$ a \omterm{def:g10:algebra:equation}{equação} diofantina
\[
17x - 40y = 1,
\]
e descreva em seguida todas as soluções de $17x - 40y = 6$.
\end{exercise}

\begin{exercise}[$\star\star$]\label{exo:g12:arith:8}
Mostre que $\sqrt2$ é irracional usando a unicidade da fatoração em \omterm{def:g12:arith:prime}{primos}
(compare o expoente de $2$ nos dois lados de $a^2 = 2b^2$).
\end{exercise}

\begin{exercise}[$\star\star\star$]\label{exo:g12:arith:9}
Seja $p$ um \omterm{def:g12:arith:prime}{primo}.
\begin{enumerate}
  \item Mostre que, para $1 \leq k \leq p - 1$, $p$ \omterm{def:g12:arith:divides}{divide}
        $\dbinom{p}{k}$. (Sugestão: use $k\binom pk = p\binom{p-1}{k-1}$,
        \cref{exo:g12:comb:7}, e o lema de Gauss.)
  \item Deduza, por indução em $a \geq 0$, outra demonstração do pequeno
        teorema de Fermat na forma $a^p \equiv a \pmod p$.
\end{enumerate}
\end{exercise}

\begin{exercise}[$\star\star\star$]\label{exo:g12:arith:10}
\emph{(Problema chinês dos restos.)} Encontre todos os inteiros $n$ tais
que
\[
n \equiv 2 \pmod 3, \qquad n \equiv 3 \pmod 5, \qquad n \equiv 2 \pmod 7 .
\]
(Sugestão: resolva as duas primeiras condições e incorpore a terceira; os
coeficientes de Bézout ajudam.)
\end{exercise}

\section{Problema: Códigos secretos e dígitos verificadores}

\begin{problem}[{Problema de fim de semana --- as congruências
protegem cada código de barras e cada cartão de crédito, e o pequeno
teorema de Fermat opera a fechadura dos segredos do mundo}]
\label{pb:g12:arith:1}
G.\,H.~Hardy gabava-se, em 1940, de que a teoria dos números era
``imaculada'' por aplicações. Oitenta anos depois, cada bipe de
código de barras, cada pagamento com cartão e cada mensagem
cifrada o contradizem --- e exatamente com as ferramentas deste
capítulo: \omterm{def:g12:arith:congruence}{congruências} (\cref{prop:g12:arith:congops}), inversos de
Bézout (\cref{thm:g12:arith:bezout}) e o pequeno teorema de Fermat
(\cref{exo:g12:arith:9}). Este problema confere os códigos, arromba
uma versão de brinquedo da fechadura e aprende por que a fechadura
verdadeira aguenta.

\medskip
\noindent\textbf{Parte I --- Fluência com \omterm{def:g12:arith:congruence}{congruências}.}
\begin{enumerate}
  \item Calcule $2026 \bmod 7$; depois o último algarismo de
        $7^{100}$ (encontre o ciclo das potências de $7$ módulo
        $10$).
  \item Exponenciação rápida (\cref{met:g12:arith:powers}):
        calcule $5^{117} \bmod 13$ (parta de
        $5^2 \equiv -1$).
  \item Resolva $3x \equiv 5 \pmod 7$.
  \item Rode o \omterm{prop:g12:arith:euclidalgo}{algoritmo de Euclides} em $(97, 35)$, substitua de
        volta para achar inteiros $u, v$ com $97u + 35v = 1$ e
        deduza o inverso de $35$ módulo $97$.
  \item Enuncie com precisão quando $a$ é invertível módulo $n$ e
        qual teorema entrega o inverso.
\end{enumerate}

\medskip
\noindent\textbf{Parte II --- Dígitos verificadores.}
\begin{enumerate}[resume]
  \item ISBN-10: os dez algarismos $d_1 \dots d_{10}$ do código
        de um livro devem satisfazer
        $10d_1 + 9d_2 + \dots + 2d_9 + 1d_{10} \equiv 0
        \pmod{11}$. Verifique o ISBN real
        $0\,306\,40615\,2$.
  \item Demonstre que o esquema ISBN detecta \emph{todo} erro em
        um único algarismo: se um algarismo muda de
        $d \not\equiv 0$, a soma ponderada muda de $w d$ com
        $1 \leq w \leq 10$ --- por que isso nunca pode ser
        $\equiv 0 \pmod{11}$ (\cref{thm:g12:arith:gauss})?
  \item Demonstre que ele também detecta toda troca de dois
        algarismos adjacentes (distintos). Depois explique o
        segredo do projeto: que propriedade de $11$ fez as duas
        demonstrações funcionarem, e o que poderia dar errado com
        o módulo $10$?
  \item Os códigos de barras EAN-13 ponderam os algarismos por
        $1, 3, 1, 3, \dots$ módulo $10$. Calcule o dígito
        verificador que completa $978\,2940199\,05$. Que trocas
        de algarismos adjacentes o EAN \emph{deixa} de detectar?
        (Quando é $2(a - b) \equiv 0 \pmod{10}$?)
  \item Os cartões de crédito usam o esquema de Luhn: da direita
        para a esquerda, dobre um algarismo sim, outro não
        (subtraindo $9$ quando o dobro passa de $9$), some tudo e
        exija um múltiplo de $10$. Verifique o número de teste
        $4539\,1488\,0343\,6467$.
  \item Em uma frase: o que o módulo primo comprou para o ISBN e
        que o EAN e o Luhn, presos ao $10$, não podem ter?
\end{enumerate}

\medskip
\noindent\textbf{Parte III --- A fechadura de Fermat.}
\begin{enumerate}[resume]
  \item Uma armadilha antes do tesouro: calcule
        $2^{10} \bmod 341$, deduza $2^{340} \bmod 341$ --- e
        depois fatore $341$. O que esse exemplo (um
        \emph{pseudoprimo de Fermat}) diz sobre usar o pequeno
        teorema de Fermat como teste de primalidade?
  \item RSA em miniatura: tome $p = 3$, $q = 11$, de modo que
        $n = 33$ e $(p-1)(q-1) = 20$; o expoente público é
        $e = 3$. Encontre o expoente privado $d$ com
        $3d \equiv 1 \pmod{20}$ (o método da questão 4).
  \item Cifre a mensagem $m = 4$: calcule
        $c = m^3 \bmod 33$.
  \item Decifre: calcule $c^d \bmod 33$ (use
        $c \equiv -2 \pmod{33}$) e recupere a mensagem.
  \item Por que a decifragem sempre funciona: mostre que
        $m^{21} \equiv m$ tanto módulo $3$ quanto módulo $11$ (o
        pequeno teorema de Fermat em cada mundo) e conclua módulo
        $33$ (o \cref{thm:g12:arith:gauss} cola as duas
        \omterm{def:g12:arith:congruence}{congruências}). Onde entrou a forma especial
        $1 + 20k$ de $21 = ed$?
  \item A segurança da fechadura: todo mundo conhece $n$ e $e$;
        recuperar $d$ exige $(p-1)(q-1)$ e, portanto, os fatores
        de $n$. Nosso $33$ se fatora à primeira vista --- por que
        o mesmo esquema, com $n$ de seiscentos algarismos,
        protege os bancos do mundo? (Uma frase sobre a assimetria
        entre multiplicar e \omterm{def:g10:algebra:expand}{fatorar}.)
\end{enumerate}

\medskip
\noindent\textbf{Parte IV --- Clássicos.}
\begin{enumerate}[resume]
  \item A antiga contagem chinesa de soldados (compare com o
        \cref{exo:g12:arith:10}): um contingente deixa resto $2$
        quando enfileirado de $3$ em $3$ e resto $3$ quando
        enfileirado de $5$ em $5$. Encontre todos os efetivos
        possíveis e explique por que a resposta é única módulo
        $15$.
  \item Demonstrações de uma linha, enfim: de $10 \equiv 1
        \pmod 9$, demonstre que todo número é congruente à soma
        de seus algarismos módulo $9$; de $10 \equiv -1
        \pmod{11}$, deduza a regra da soma alternada para o $11$.
        (O volume do ensino fundamental demonstrou isso com
        álgebra explícita --- admire a compressão.)
  \item Final --- Hardy contra o código de barras: recapitule a
        caixa de ferramentas do capítulo (aritmética das
        \omterm{def:g12:arith:congruence}{congruências}, inversos de Bézout, pequeno teorema de
        Fermat, colagem de módulos primos entre si) e onde cada
        uma se encaixou neste problema; depois dê o veredicto
        moderno sobre o ``imaculada''.
\end{enumerate}
\end{problem}
